\section{Comparing finite-size centers with proper clauses}
\label{sec:direct-comparison}

The variance estimate controls fluctuations at one size.  We now compare
the corresponding centers at different sizes, using proper clauses at
every step.  The comparison keeps the weights of the two variable blocks;
this is what makes its error uniform even when one block is much larger
than the other.  Once these probability comparisons are established, the
deterministic balanced-comparison lemma of the threshold companion turns
them into convergence. This gives a second route to convergence from a
concentration estimate, without using the companion's limiting-threshold
theorem as Section~\ref{sec:interfaces} does. The theorem below takes
concentration as an input, so it also applies to the independent,
weaker variance proof in Appendix~\ref{sec:appendix-c}.

Interpolation between a system and independent parts appears in the
thermodynamic-limit work of Guerra and Toninelli~\cite{GuerraToninelli}
and the diluted-system work of Franz and Leone~\cite{FranzLeone} and
Franz, Leone, and Toninelli~\cite{FranzLeoneToninelli}.
Those works concern free energies under their respective model
hypotheses; they provide methodological context for the size comparison.
The one-clause interpolation here uses frozen variables, as in the
satisfiability interpolation of Bayati, Gamarnik, and Tetali
\cite[Section~4, Proposition~2]{BGT}.  Their clause law samples variable
tuples with replacement.  Here the clauses have distinct variables, so
the exact killing probabilities contain falling factorials.
We control the resulting error locally before summing it over the trials.

\subsection{A center theorem with an explicit concentration input}

For fixed $k\ge3$ and $n\ge k$, let $H_n$ be the index of the first
unsatisfiable prefix of independent uniform proper $k$-clauses on $n$
variables, and write $P_n(j)=\PP(H_n>j)$ for integers $j\ge0$.
Fix a real $B>0$ and set
\[
 L_n=\lfloor Bn\rfloor,\qquad T_n=\min\{H_n,L_n\},\qquad
 \mu_n=\EE T_n,\qquad e_n=\mu_n/n.
\]
The endpoint for this first-unsatisfiable convention is strict:
\begin{equation}\label{cmp:prefix-endpoint}
 P_n(j)=\PP(T_n>j)\qquad(0\le j<L_n).
\end{equation}
At $j=L_n$, the probability on the right is zero, whereas $P_n(L_n)$
need not vanish.  All statements below use only indices strictly below
the cap when applying~\eqref{cmp:prefix-endpoint}.

\begin{theorem}[Proper-clause center comparison from concentration]
\label{thm:proper-center-module}\label{cmp:modular}
\label{cmp:centers}\label{cmp:convergence}
Fix $k\ge3$, $B>0$, and $0<\delta<1/2$.  Suppose that there are constants
$0<a\le A<B$ and numbers $\eta_n\ge0$ tending to zero such that, for all
sufficiently large $n$,
\begin{equation}\label{cmp:module-input}
 a\le e_n\le A,\qquad
 \PP\bigl(|T_n-\mu_n|\ge n^{1-\delta}\bigr)\le\eta_n.
\end{equation}
Then there is a finite constant $K\ge0$ such that, for all sufficiently large
$n=\min\{n_1,n_2\}$, with no restriction on $n_1/n_2$,
\begin{equation}\label{cmp:sum}
 e_{n_1+n_2}\ge\min\{e_{n_1},e_{n_2}\}-K n^{-\delta},
\end{equation}
and, for all sufficiently large $n$,
\begin{equation}\label{cmp:neighbor}
 e_{n+1}\ge e_n-K n^{-\delta}.
\end{equation}
Moreover $e_n$ converges to a number $\alpha\in[a,A]$, and
\begin{equation}\label{cmp:strict-sides}
 P_n(\lfloor cn\rfloor)\longrightarrow
 \begin{cases}1,&0\le c<\alpha,\\0,&c>\alpha.\end{cases}
\end{equation}
For every sufficiently large $s$ and every $N\ge s^2$, one has the
one-sided estimate
\begin{equation}\label{cmp:tree-bound}
 e_N\ge e_s-K(1+B_\delta)s^{-\delta},\qquad
 B_\delta=\frac{(3/2)^\delta}{1-(2/3)^\delta}.
\end{equation}
\end{theorem}

The theorem allows an arbitrary vanishing concentration error; it does
not require a specified polynomial rate.  In particular, Chebyshev's
inequality supplies~\eqref{cmp:module-input} whenever
\begin{equation}\label{cmp:variance-input}
 \Var(T_n)=o\bigl(n^{2-2\delta}\bigr).
\end{equation}
The proof uses concentration twice: to make the two local formulas likely
to be satisfiable, and to turn a likely satisfiable diluted formula back
into a lower bound for the center at the combined size.

\subsection{Local clauses and diluted global clauses}

\begin{lemma}[Proper-clause trial comparison]\label{cmp:trial-comparison}
Partition $N=n_1+n_2$ variables into blocks of sizes $n_1,n_2\ge k$,
and put $w_j=n_j/N$ and $n=\min\{n_1,n_2\}$.  Fix an integer $M\ge0$
and $0<\varepsilon<1$.  Each of $M$ independent local trials chooses
block $j$ with probability $w_j$ and inserts an independent uniform
proper $k$-clause in that block.  Each of $M$ independent diluted global
trials inserts an independent uniform proper $k$-clause on all $N$
variables with probability $1-\varepsilon$, and inserts no clause
otherwise.  Let $Q_{\mathrm{loc}}$ and $Q_{\mathrm{dil}}$ be the
respective probabilities that the resulting formulas are satisfiable.
Then
\begin{align}
 Q_{\mathrm{dil}}
 &\ge Q_{\mathrm{loc}}-
 C_k M\varepsilon^{-k}\sum_{j=1}^{2}\frac{w_j}{n_j^k},
 & C_k&=\frac{[k(k-1)]^k}{2^k},\label{cmp:weighted-replacement}\\
 Q_{\mathrm{dil}}
 &\ge Q_{\mathrm{loc}}-\frac{C_kM}{\varepsilon^k n^k}.
 \label{cmp:coarse-replacement}
\end{align}
For $N=n+1$, if all local trials instead use a fixed main set of
$n\ge k$ variables, then
\begin{equation}\label{cmp:neighbor-replacement}
 Q_{\mathrm{dil}}\ge Q_{\mathrm{loc}}
       -\frac{D_kM}{\varepsilon^k n^k},\qquad
 D_k=\frac{k^{2k}}{2^k}.
\end{equation}
These statements impose no restriction on a block-size ratio.
\end{lemma}

\begin{proof}
Condition on every trial except the one being replaced.  If the residual
formula is unsatisfiable, both conditional satisfiability probabilities
are zero.  Otherwise let $S$ be its nonempty solution set, and call a
variable frozen when it has the same value in every member of $S$.
Among any specified $v\ge k$ variables of which $b$ are frozen, a proper
$k$-clause kills $S$ with probability
\[
 h_k(b,v)=\frac{(b)_k}{2^k(v)_k}.
\]
Indeed, each of the selected variables must be frozen and each sign must
oppose its frozen value.  With $x=b/v$ and $y_+=\max\{y,0\}$,
\begin{equation}\label{cmp:factorial-bounds}
 \frac{(x-(k-1)/v)_+^k}{2^k}
 \le h_k(b,v)\le\frac{x^k}{2^k}.
\end{equation}
For $b\ge k$, each factor $(b-r)/(v-r)$, $0\le r<k$, is at most $b/v$
and at least $(b-k+1)/v$.  For $b<k$, both the middle expression and
the positive part on the left vanish.  This proves the bounds.

Let $\beta_j$ be the frozen fraction in block $j$, always with respect
to the entire residual set $S$; its defining formula may contain clauses
crossing the two blocks.  By convexity and~\eqref{cmp:factorial-bounds},
the killing probability of a nondiluted global clause is at most
\[
 2^{-k}\left(\sum_jw_j\beta_j\right)^k
 \le 2^{-k}\sum_jw_j\beta_j^k,
\]
whereas the local killing probability is at least
$2^{-k}\sum_jw_j(\beta_j-(k-1)/n_j)_+^k$.
If $\beta_j\ge k(k-1)/(\varepsilon n_j)$, Bernoulli's inequality gives
\[
 (\beta_j-(k-1)/n_j)_+^k
 \ge (1-\varepsilon/k)^k\beta_j^k
 \ge (1-\varepsilon)\beta_j^k.
\]
Below that threshold,
$\beta_j^k\le[k(k-1)]^k/(\varepsilon^k n_j^k)$.
It follows that the diluted global killing probability exceeds the local
one by at most $C_k\varepsilon^{-k}\sum_jw_j/n_j^k$.
The estimate holds for every residual formula.  Replacing the trials one
at a time, conditioning on all other independent trials at each step,
and then averaging and telescoping proves~\eqref{cmp:weighted-replacement}.
Since $n_j\ge n$ and $w_1+w_2=1$, it also proves the coarser bound.

For neighboring sizes, let $b$ count the frozen variables in the main
$n$-set and let $x\in\{0,1\}$ indicate whether the extra variable is
frozen.  The global frozen fraction $\gamma=(b+x)/(n+1)$ obeys
$b/n\ge\gamma-1/(n+1)$.  The local killing probability is therefore
at least $2^{-k}(\gamma-k/n)_+^k$, while the nondiluted global one is at
most $2^{-k}\gamma^k$.  Splitting at
$\gamma=k^2/(\varepsilon n)$ and applying the same Bernoulli inequality
bounds the per-trial loss by $D_k/(\varepsilon^k n^k)$.  Telescoping gives
\eqref{cmp:neighbor-replacement}.
\end{proof}

At $k=3$ the constants are exactly $C_3=27$ and $D_3=729/8$.
The stronger weighted bound will be used for unequal blocks.  The coarser
bound remains useful when the block sizes are known to be comparable.

\subsection{From a random clause count to a center}

We first record the integer step needed in both center comparisons.
From~\eqref{cmp:module-input} and~\eqref{cmp:prefix-endpoint}, for every
integer $0\le j<L_n$ and all sufficiently large $n$,
\begin{equation}\label{cmp:concentration}
 \begin{aligned}
 j\le\mu_n-n^{1-\delta}&\quad\Longrightarrow\quad P_n(j)\ge1-\eta_n,\\
 j\ge\mu_n+n^{1-\delta}&\quad\Longrightarrow\quad P_n(j)\le\eta_n.
 \end{aligned}
\end{equation}
Because $\eta_n\to0$, its tail supremum
$\overline\eta_n=\sup_{v\ge n}\eta_v$ also tends to zero.  This will
provide uniformity when one block is arbitrarily larger than the other.

Suppose, for a fixed $a_0>0$, that
$a_0N\le M\le AN$, $0<\varepsilon\le1/2$, and
$Q_{\mathrm{dil}}\ge1-\zeta$.  The number of inserted clauses is
$J\sim\operatorname{Bin}(M,1-\varepsilon)$.  Conditional on $J$, the
inserted clauses have exactly the independent proper-clause law, so
$Q_{\mathrm{dil}}=\EE P_N(J)$.  Put
\[
 j=\left\lfloor(1-\varepsilon)M-N^{1-\delta}\right\rfloor.
\]
For all sufficiently large $N$, the lower bound on $M$ gives $j\ge0$,
and $j\le M\le AN<L_N$ because $A<B$.  Also $\Var(J)\le AN$, whence
\[
 \PP(J<j)\le A N^{-1+2\delta},\qquad
 Q_{\mathrm{dil}}\le P_N(j)+\PP(J<j).
\]
The second inequality uses monotonicity of $P_N$.  If
$\zeta+AN^{-1+2\delta}+\eta_N<1$, the upper implication
in~\eqref{cmp:concentration} forces $j<\mu_N+N^{1-\delta}$.
The floor in the definition of $j$ then gives
\begin{equation}\label{cmp:count-transfer}
 e_N>\frac{(1-\varepsilon)M}{N}-2N^{-\delta}-N^{-1}
 \ge\frac{(1-\varepsilon)M}{N}-3N^{-\delta}.
\end{equation}
We will apply this implication with all three errors tending to zero.

\begin{proof}[Proof of Theorem~\ref{thm:proper-center-module}]
For the two-block comparison, put
$N=n_1+n_2$, $n=\min\{n_1,n_2\}$,
$e_* =\min\{e_{n_1},e_{n_2}\}$, and choose
\[
 \varepsilon=n^{-\delta},\qquad
 M=\left\lfloor(e_*-2\varepsilon)N\right\rfloor.
\]
For sufficiently large $n$, the anchors imply
$(a/4)N\le M\le AN$ and $\varepsilon\le1/2$.
In the local model let $X_j$ count the clauses in block $j$.  Its marginal
law is binomial, and
\[
 \EE X_j=Mn_j/N\le\mu_{n_j}-2\varepsilon n_j
       \le\mu_{n_j}-2n_j^{1-\delta},\qquad
 \Var(X_j)\le A n_j.
\]
Thus
$\PP(X_j>\mu_{n_j}-n_j^{1-\delta})\le A n_j^{-1+2\delta}$.
Conditional on the block choices, the clauses inside the two blocks have
independent proper laws.  On the event
$X_j\le\mu_{n_j}-n_j^{1-\delta}$, the integer count lies below $L_{n_j}$
for large $n_j$, because $\mu_{n_j}\le A n_j$ and $A<B$.
The lower implication in~\eqref{cmp:concentration} and a union bound
therefore give
\[
 Q_{\mathrm{loc}}\ge1-
 \sum_{j=1}^{2}\bigl(A n_j^{-1+2\delta}+\eta_{n_j}\bigr).
\]
The right-hand error is at most
$2A n^{-1+2\delta}+2\overline\eta_n$, which tends to zero uniformly
over all pairs of sizes with minimum $n$.

The weighted replacement loss has the same uniformity:
\begin{align}
 C_kM\varepsilon^{-k}\sum_j\frac{w_j}{n_j^k}
 &=C_k\frac{M}{N}\varepsilon^{-k}\sum_jn_j^{1-k}\notag\\
 &\le2AC_k n^{1-k+k\delta}=o(1).
 \label{cmp:uniform-replacement-loss}
\end{align}
Here $1-k+k\delta<0$ for $k\ge3$ and $\delta<1/2$.
The factor $N$ has canceled, so no upper bound on $N/n$ is used.
Lemma~\ref{cmp:trial-comparison} therefore gives
$Q_{\mathrm{dil}}\ge1-\zeta_n$, where $\zeta_n\to0$ uniformly over
these pairs.  Since $N\ge2n$, the count error at size $N$ and
$\eta_N\le\overline\eta_n$ also vanish uniformly.  Applying
\eqref{cmp:count-transfer}, and keeping the floor in $M$, yields
\[
 e_N\ge(1-\varepsilon)(e_*-2\varepsilon-N^{-1})-3N^{-\delta}
       \ge e_*-K n^{-\delta},
\]
where the last inequality uses $e_*\le A$, $N\ge n$, and $\delta<1$.
This proves~\eqref{cmp:sum} for arbitrary size ratios.

For the neighboring comparison take $N=n+1$, put
$\varepsilon=n^{-\delta}$, and let
$M=\lfloor\mu_n-2n^{1-\delta}\rfloor$ local trials all use the main
$n$ variables.  Again $(a/4)N\le M\le AN$ for sufficiently large $n$.
The local formula has exactly the law at size $n$ and count $M$, so
$Q_{\mathrm{loc}}\ge1-\eta_n$ by~\eqref{cmp:concentration}.
The neighboring replacement loss is
\[
 \frac{D_kM}{\varepsilon^k n^k}=O_{k,A}(n^{1-k+k\delta})=o(1).
\]
Equation~\eqref{cmp:count-transfer} now gives
\[
 e_{n+1}\ge(1-\varepsilon)
       \frac{ne_n-2n^{1-\delta}-1}{n+1}-3(n+1)^{-\delta}
       \ge e_n-Kn^{-\delta}.
\]
For the last step, use $e_n\le A$ and
$1-(1-\varepsilon)n/(n+1)\le\varepsilon+1/(n+1)$.
This proves~\eqref{cmp:neighbor}.

It remains to explain the deterministic step and check its inputs.
The geometric summation of errors parallels Abbe and Montanari's
approximate sum-superadditivity argument
\cite[arXiv version, Lemma~8 and Remark~3]{AM}.  The required minimum
comparison is the threshold companion's Lemma~\ThresholdBalanced{}
(``Balanced comparisons'')~\cite{FixedKThreshold}.  For an integer
$n_0\ge1$, it states that if a bounded real sequence $(u_n)_{n\ge n_0}$
satisfies
\[
 u_{r+s}\ge\min\{u_r,u_s\}-K\min(r,s)^{-\delta}
 \quad\text{for }1/3\le r/s\le3,\qquad
 u_{s+1}\ge u_s-Ks^{-\delta},
\]
for all $r,s\ge n_0$ in the indicated ratio range and all $s\ge n_0$,
respectively, with fixed $K\ge0$ and $\delta>0$, then it converges and,
for every $s\ge\max\{n_0,2\}$ and every $N\ge s^2$,
$u_N\ge u_s-K(1+B_\delta)s^{-\delta}$ with $B_\delta$ as above.
The lemma is a deterministic minimum comparison: it uses leaves of sizes
$s,s+1$, and its binary tree has child-size ratios at most three.
Choose $n_0$ beyond the cutoffs in our anchors and two comparisons.
Then the boundedness in~\eqref{cmp:module-input},
the comparable-size part of \eqref{cmp:sum}, and~\eqref{cmp:neighbor}
verify every hypothesis with $u_n=e_n$. The lemma gives~\eqref{cmp:tree-bound} and convergence to
some $\alpha\in[a,A]$.

For $0\le c<\alpha$, eventually
$\lfloor cn\rfloor\le\mu_n-n^{1-\delta}$ and $\lfloor cn\rfloor<L_n$.
For $\alpha<c<B$, eventually
$\mu_n+n^{1-\delta}\le\lfloor cn\rfloor<L_n$.
In both cases the fixed positive density gap eventually exceeds
$n^{1-\delta}+1$, accounting for the floor.  Apply
\eqref{cmp:concentration}.  Since $\alpha\le A<B$, choose one fixed
$d\in(A,B)$; monotonicity extends the zero limit from $d$ to all
$c\ge B$.  This proves~\eqref{cmp:strict-sides}.
\end{proof}

The comparison~\eqref{cmp:tree-bound} is one-sided.  It proves
convergence of the finite-size means, but gives no two-sided estimate for
$\mu_n-\alpha n$ on the fluctuation scale of the variance theorem.

\subsection{Checking the inputs for the variance estimates}

For the main application put $B_k=2^{k+1}$ and take $B=B_k$; at $k=3$
we may also take $B=10$.  We verify the anchors without assuming threshold
convergence.  Let $q_k=1-2^{-k}$.  The assignment first moment gives
$P_n(j)\le2^nq_k^j$, hence, for every positive integer $d$,
\begin{equation}\label{cmp:mean-upper}
 \EE H_n=\sum_{j=0}^{\infty}P_n(j)
       \le dn+2^k(2q_k^d)^n.
\end{equation}
For $d=2^k$, one has $2q_k^{2^k}<2/e<1$.  For $k=3$, the sharper choice
$d=6$ satisfies $2(7/8)^6<1$.  Since $T_n\le H_n$, it follows that
eventually
\begin{equation}\label{cmp:center-upper}
 e_n\le A_k<B,
 \qquad A_k=\begin{cases}7,&k=3,\\2^k+1,&k\ge4.\end{cases}
\end{equation}

For the positive anchor, fix $0<c<\min\{e^{-2},B\}$ and let
$j=\lfloor cn\rfloor$.  Lemma~\ref{lem:sparse-satisfiability} gives
$P_n(j)=1-O_c(n^{-2})$ for the same independent proper-clause model.
Since $c<B$, one has $j<L_n$ for sufficiently large $n$.  Therefore
$\mu_n\ge jP_n(j)=cn-o(n)$, and with $c_0=c/2$ we obtain the full anchors
\begin{equation}\label{cmp:center-bounds}
 c_0\le e_n\le A_k<B
 \qquad\text{for all sufficiently large }n.
\end{equation}

The capped variance theorem gives
\[
 \frac{\Var(T_n)}{n^{2-2\delta}}
 =\begin{cases}
 O(n^{-1+2\delta}\log n),&k=3,\\
 O_k(n^{-1+2\delta}),&k\ge4,
 \end{cases}
\]
which tends to zero for every fixed $0<\delta<1/2$.
Theorem~\ref{thm:proper-center-module} now applies with $a=c_0$ and
$A=A_k$.  Its strict-side limit is the same $\alpha_k$ as in the threshold
companion~\cite[Theorem~\ThresholdTheorem]{FixedKThreshold}: if the two
limits differed, a density strictly between them would be assigned opposite
limiting probabilities.  In particular, the proper-clause comparisons
\eqref{cmp:sum}--\eqref{cmp:neighbor} hold for every such fixed $\delta$.

The weaker polynomial variance estimate of the threshold companion
\cite[Proposition~\ThresholdVariance]{FixedKThreshold} can be used as a
separate input.  For $a=B_kn$ and the same last-SAT cap $b=a$, the exact
identity in Proposition~\ref{prop:cap-bridge} is
\[
 \min\{H_n-1,a\}=\min\{H_n,a\}-1+\ind{H_n>a}.
\]
That proposition gives
$\Var(\min\{H_n,a\})\le\Var(\min\{H_n-1,a\})$.
Consequently the companion's $O_k(n^{1+2/k})$ bound verifies
\eqref{cmp:variance-input} for every
$0<\delta<(k-2)/(2k)$, including the companion's stated choice
$\delta=(k-2)/(4k)$.  This substitution uses the companion's own
variance estimate.  The range $\delta<1/2$ for the local variance bounds
is also strict; no endpoint assertion is made.
